% Chapter 0 (front matter): Acknowledgments
%
% Generated 2026-09-08 (Round-9 / session continuation) to fill the
% missing acknowledgments that was flagged by the Tier-6 audit 2026-09-07.
%
% Cross-referenced from thesis/main.tex % Chapter 0 (front matter): Acknowledgments
%
% Generated 2026-09-08 (Round-9 / session continuation) to fill the
% missing acknowledgments that was flagged by the Tier-6 audit 2026-09-07.
%
% Cross-referenced from thesis/main.tex % Chapter 0 (front matter): Acknowledgments
%
% Generated 2026-09-08 (Round-9 / session continuation) to fill the
% missing acknowledgments that was flagged by the Tier-6 audit 2026-09-07.
%
% Cross-referenced from thesis/main.tex % Chapter 0 (front matter): Acknowledgments
%
% Generated 2026-09-08 (Round-9 / session continuation) to fill the
% missing acknowledgments that was flagged by the Tier-6 audit 2026-09-07.
%
% Cross-referenced from thesis/main.tex \input{chapters/00_acknowledgments}.
% Wordcount target: ~400 words.
%
% Per docs/decisions/OPEN-QUESTIONS-FOR-HUMAN-2026-09-07.md Q4 (Option E):
% communities are acknowledged in the aggregate, not named individually,
% until FPIC engagement is complete.

\chapter*{Acknowledgments}
\addcontentsline{toc}{chapter}{Acknowledgments}

This thesis is the product of five years of part-time research
conducted while the author completed the \textit{Facultad de
Ciencias Agrarias} (FADA) undergraduate program at the
\textit{Universidad Nacional de Asunción} (UNA), with primary
funding from the author's personal stipend and infrastructure
support from FADA-UNA computing resources. No external grant
funding supported the work.

The author thanks the \textbf{ten indigenous communities} of the
Gran Chaco whose territories are analyzed in
Chapter~\ref{ch:p0012-yvy}. The deforestation disparity documented
in this thesis (3.0$\times$ relative to the national rate) is
reported with anonymized per-territory data per the CARE
Principles for Indigenous Data Governance; per-community
disclosure is contingent on the Free, Prior, and Informed Consent
(FPIC) engagement process documented in
Appendix~\ref{ch:appendix:ethics}. The communities' stewardship of
the Chaco's forests is the empirical baseline against which this
thesis measures structural injustice.

The author thanks the institutional partners who provided data,
validation, or methodological guidance under the formal
partnership-letter process documented at
\texttt{docs/partnerships/ONE-PAGERS-2026-09-07.md}:

\begin{itemize}
    \item \textbf{INFONA} (\textit{Instituto Forestal Nacional}) —
      institutional partner for P0011 Yvutu (Chaco deforestation
      detection); pending partnership letter.
    \item \textbf{INDI} (\textit{Instituto Paraguayo del Indígena})
      — institutional partner for P0012 Yvy (indigenous territory
      analysis); pending FPIC engagement per ILO~169.
    \item \textbf{SENEPA} (\textit{Servicio Nacional de Erradicación
      del Paludismo}) — public-health context for P0035 Tatakua
      (air-quality forecasting).
    \item \textbf{Verra} (\textit{Verified Carbon Standard}) —
      carbon-market registry data for P0010 Yvyra.
    \item \textbf{Guyra Paraguay} — wildlife camera-trap imagery
      for P0026 Kai.
    \item \textbf{WWF Paraguay} — Chaco monitoring reports (cited
      in the literature review).
    \item \textbf{INBIO} (\textit{Instituto de Biotecnología
      Agrícola}) — soybean-yield ground truth for P0025 Yrupe.
    \item \textbf{OpenAQ} — open air-quality station network for
      P0035 Tatakua.
\end{itemize}

The author thanks \textbf{Juan Carlos Cristaldo} (FADA-UNA,
\textit{Facultad de Ciencias Exactas y Naturales}) for thesis
adviser engagement, the \textbf{FADA-UNA thesis committee} for
institutional review, and the open-source community whose tools
(Hansen GFC, MapBiomas Paraguay, OpenStreetMap, Prithvi, AlphaEarth,
DINOv2, YOLOv8, TimesFM, the Python scientific stack) made this
thesis possible.

Finally, the author thanks the \textbf{IEEE Geoscience and Remote
Sensing Society}, the \textbf{Open Data Institute}, and the
\textbf{Free Software Foundation} for sustaining the open-science
infrastructure that this thesis depends on. Without their
commitment to reproducible science and rights-aware data governance,
none of the findings here would be publishable.

\textbf{Author:} Iván Weiss Van der Pol \\
FADA-UNA, San Lorenzo, Paraguay \\
September 2026
.
% Wordcount target: ~400 words.
%
% Per docs/decisions/OPEN-QUESTIONS-FOR-HUMAN-2026-09-07.md Q4 (Option E):
% communities are acknowledged in the aggregate, not named individually,
% until FPIC engagement is complete.

\chapter*{Acknowledgments}
\addcontentsline{toc}{chapter}{Acknowledgments}

This thesis is the product of five years of part-time research
conducted while the author completed the \textit{Facultad de
Ciencias Agrarias} (FADA) undergraduate program at the
\textit{Universidad Nacional de Asunción} (UNA), with primary
funding from the author's personal stipend and infrastructure
support from FADA-UNA computing resources. No external grant
funding supported the work.

The author thanks the \textbf{ten indigenous communities} of the
Gran Chaco whose territories are analyzed in
Chapter~\ref{ch:p0012-yvy}. The deforestation disparity documented
in this thesis (3.0$\times$ relative to the national rate) is
reported with anonymized per-territory data per the CARE
Principles for Indigenous Data Governance; per-community
disclosure is contingent on the Free, Prior, and Informed Consent
(FPIC) engagement process documented in
Appendix~\ref{ch:appendix:ethics}. The communities' stewardship of
the Chaco's forests is the empirical baseline against which this
thesis measures structural injustice.

The author thanks the institutional partners who provided data,
validation, or methodological guidance under the formal
partnership-letter process documented at
\texttt{docs/partnerships/ONE-PAGERS-2026-09-07.md}:

\begin{itemize}
    \item \textbf{INFONA} (\textit{Instituto Forestal Nacional}) —
      institutional partner for P0011 Yvutu (Chaco deforestation
      detection); pending partnership letter.
    \item \textbf{INDI} (\textit{Instituto Paraguayo del Indígena})
      — institutional partner for P0012 Yvy (indigenous territory
      analysis); pending FPIC engagement per ILO~169.
    \item \textbf{SENEPA} (\textit{Servicio Nacional de Erradicación
      del Paludismo}) — public-health context for P0035 Tatakua
      (air-quality forecasting).
    \item \textbf{Verra} (\textit{Verified Carbon Standard}) —
      carbon-market registry data for P0010 Yvyra.
    \item \textbf{Guyra Paraguay} — wildlife camera-trap imagery
      for P0026 Kai.
    \item \textbf{WWF Paraguay} — Chaco monitoring reports (cited
      in the literature review).
    \item \textbf{INBIO} (\textit{Instituto de Biotecnología
      Agrícola}) — soybean-yield ground truth for P0025 Yrupe.
    \item \textbf{OpenAQ} — open air-quality station network for
      P0035 Tatakua.
\end{itemize}

The author thanks \textbf{Juan Carlos Cristaldo} (FADA-UNA,
\textit{Facultad de Ciencias Exactas y Naturales}) for thesis
adviser engagement, the \textbf{FADA-UNA thesis committee} for
institutional review, and the open-source community whose tools
(Hansen GFC, MapBiomas Paraguay, OpenStreetMap, Prithvi, AlphaEarth,
DINOv2, YOLOv8, TimesFM, the Python scientific stack) made this
thesis possible.

Finally, the author thanks the \textbf{IEEE Geoscience and Remote
Sensing Society}, the \textbf{Open Data Institute}, and the
\textbf{Free Software Foundation} for sustaining the open-science
infrastructure that this thesis depends on. Without their
commitment to reproducible science and rights-aware data governance,
none of the findings here would be publishable.

\textbf{Author:} Iván Weiss Van der Pol \\
FADA-UNA, San Lorenzo, Paraguay \\
September 2026
.
% Wordcount target: ~400 words.
%
% Per docs/decisions/OPEN-QUESTIONS-FOR-HUMAN-2026-09-07.md Q4 (Option E):
% communities are acknowledged in the aggregate, not named individually,
% until FPIC engagement is complete.

\chapter*{Acknowledgments}
\addcontentsline{toc}{chapter}{Acknowledgments}

This thesis is the product of five years of part-time research
conducted while the author completed the \textit{Facultad de
Ciencias Agrarias} (FADA) undergraduate program at the
\textit{Universidad Nacional de Asunción} (UNA), with primary
funding from the author's personal stipend and infrastructure
support from FADA-UNA computing resources. No external grant
funding supported the work.

The author thanks the \textbf{ten indigenous communities} of the
Gran Chaco whose territories are analyzed in
Chapter~\ref{ch:p0012-yvy}. The deforestation disparity documented
in this thesis (3.0$\times$ relative to the national rate) is
reported with anonymized per-territory data per the CARE
Principles for Indigenous Data Governance; per-community
disclosure is contingent on the Free, Prior, and Informed Consent
(FPIC) engagement process documented in
Appendix~\ref{ch:appendix:ethics}. The communities' stewardship of
the Chaco's forests is the empirical baseline against which this
thesis measures structural injustice.

The author thanks the institutional partners who provided data,
validation, or methodological guidance under the formal
partnership-letter process documented at
\texttt{docs/partnerships/ONE-PAGERS-2026-09-07.md}:

\begin{itemize}
    \item \textbf{INFONA} (\textit{Instituto Forestal Nacional}) —
      institutional partner for P0011 Yvutu (Chaco deforestation
      detection); pending partnership letter.
    \item \textbf{INDI} (\textit{Instituto Paraguayo del Indígena})
      — institutional partner for P0012 Yvy (indigenous territory
      analysis); pending FPIC engagement per ILO~169.
    \item \textbf{SENEPA} (\textit{Servicio Nacional de Erradicación
      del Paludismo}) — public-health context for P0035 Tatakua
      (air-quality forecasting).
    \item \textbf{Verra} (\textit{Verified Carbon Standard}) —
      carbon-market registry data for P0010 Yvyra.
    \item \textbf{Guyra Paraguay} — wildlife camera-trap imagery
      for P0026 Kai.
    \item \textbf{WWF Paraguay} — Chaco monitoring reports (cited
      in the literature review).
    \item \textbf{INBIO} (\textit{Instituto de Biotecnología
      Agrícola}) — soybean-yield ground truth for P0025 Yrupe.
    \item \textbf{OpenAQ} — open air-quality station network for
      P0035 Tatakua.
\end{itemize}

The author thanks \textbf{Juan Carlos Cristaldo} (FADA-UNA,
\textit{Facultad de Ciencias Exactas y Naturales}) for thesis
adviser engagement, the \textbf{FADA-UNA thesis committee} for
institutional review, and the open-source community whose tools
(Hansen GFC, MapBiomas Paraguay, OpenStreetMap, Prithvi, AlphaEarth,
DINOv2, YOLOv8, TimesFM, the Python scientific stack) made this
thesis possible.

Finally, the author thanks the \textbf{IEEE Geoscience and Remote
Sensing Society}, the \textbf{Open Data Institute}, and the
\textbf{Free Software Foundation} for sustaining the open-science
infrastructure that this thesis depends on. Without their
commitment to reproducible science and rights-aware data governance,
none of the findings here would be publishable.

\textbf{Author:} Iván Weiss Van der Pol \\
FADA-UNA, San Lorenzo, Paraguay \\
September 2026
.
% Wordcount target: ~400 words.
%
% Per docs/decisions/OPEN-QUESTIONS-FOR-HUMAN-2026-09-07.md Q4 (Option E):
% communities are acknowledged in the aggregate, not named individually,
% until FPIC engagement is complete.

\chapter*{Acknowledgments}
\addcontentsline{toc}{chapter}{Acknowledgments}

This thesis is the product of five years of part-time research
conducted while the author completed the \textit{Facultad de
Ciencias Agrarias} (FADA) undergraduate program at the
\textit{Universidad Nacional de Asunción} (UNA), with primary
funding from the author's personal stipend and infrastructure
support from FADA-UNA computing resources. No external grant
funding supported the work.

The author thanks the \textbf{ten indigenous communities} of the
Gran Chaco whose territories are analyzed in
Chapter~\ref{ch:p0012-yvy}. The deforestation disparity documented
in this thesis (3.0$\times$ relative to the national rate) is
reported with anonymized per-territory data per the CARE
Principles for Indigenous Data Governance; per-community
disclosure is contingent on the Free, Prior, and Informed Consent
(FPIC) engagement process documented in
Appendix~\ref{ch:appendix:ethics}. The communities' stewardship of
the Chaco's forests is the empirical baseline against which this
thesis measures structural injustice.

The author thanks the institutional partners who provided data,
validation, or methodological guidance under the formal
partnership-letter process documented at
\texttt{docs/partnerships/ONE-PAGERS-2026-09-07.md}:

\begin{itemize}
    \item \textbf{INFONA} (\textit{Instituto Forestal Nacional}) —
      institutional partner for P0011 Yvutu (Chaco deforestation
      detection); pending partnership letter.
    \item \textbf{INDI} (\textit{Instituto Paraguayo del Indígena})
      — institutional partner for P0012 Yvy (indigenous territory
      analysis); pending FPIC engagement per ILO~169.
    \item \textbf{SENEPA} (\textit{Servicio Nacional de Erradicación
      del Paludismo}) — public-health context for P0035 Tatakua
      (air-quality forecasting).
    \item \textbf{Verra} (\textit{Verified Carbon Standard}) —
      carbon-market registry data for P0010 Yvyra.
    \item \textbf{Guyra Paraguay} — wildlife camera-trap imagery
      for P0026 Kai.
    \item \textbf{WWF Paraguay} — Chaco monitoring reports (cited
      in the literature review).
    \item \textbf{INBIO} (\textit{Instituto de Biotecnología
      Agrícola}) — soybean-yield ground truth for P0025 Yrupe.
    \item \textbf{OpenAQ} — open air-quality station network for
      P0035 Tatakua.
\end{itemize}

The author thanks \textbf{Juan Carlos Cristaldo} (FADA-UNA,
\textit{Facultad de Ciencias Exactas y Naturales}) for thesis
adviser engagement, the \textbf{FADA-UNA thesis committee} for
institutional review, and the open-source community whose tools
(Hansen GFC, MapBiomas Paraguay, OpenStreetMap, Prithvi, AlphaEarth,
DINOv2, YOLOv8, TimesFM, the Python scientific stack) made this
thesis possible.

Finally, the author thanks the \textbf{IEEE Geoscience and Remote
Sensing Society}, the \textbf{Open Data Institute}, and the
\textbf{Free Software Foundation} for sustaining the open-science
infrastructure that this thesis depends on. Without their
commitment to reproducible science and rights-aware data governance,
none of the findings here would be publishable.

\textbf{Author:} Iván Weiss Van der Pol \\
FADA-UNA, San Lorenzo, Paraguay \\
September 2026
